\documentclass[10pt,a4paper,landscape]{article}

\usepackage[utf8]{inputenc}
\usepackage[T1]{fontenc}
\usepackage[landscape,margin=1.4cm,top=1.2cm,bottom=1.2cm]{geometry}
\usepackage{amsmath,amssymb}
\usepackage{booktabs}
\usepackage{array}
\usepackage[table,dvipsnames]{xcolor}
\usepackage{colortbl}
\usepackage[colorlinks=true,linkcolor=Navy,citecolor=Navy,urlcolor=Navy]{hyperref}
\usepackage{caption}
\usepackage{enumitem}
\usepackage{tcolorbox}
\usepackage{multicol}

\definecolor{Navy}{RGB}{27,73,101}
\definecolor{Rust}{RGB}{169,68,66}
\definecolor{Teal}{RGB}{77,144,142}
\definecolor{Gold}{RGB}{224,159,62}
\definecolor{StrandA}{RGB}{236,241,244}
\definecolor{StrandB}{RGB}{247,240,235}
\definecolor{StrandC}{RGB}{237,244,243}
\definecolor{Ours}{RGB}{252,244,224}

\newcommand{\yes}{\textcolor{Teal}{\textbf{\checkmark}}}
\newcommand{\no}{\textcolor{black!30}{\textbf{--}}}
\newcommand{\nokey}{\textcolor{Rust}{\textbf{--}}}
\newcommand{\yeskey}{\textcolor{Rust}{\textbf{\checkmark}}}
\newcommand{\sub}[1]{{\footnotesize\textcolor{black!55}{#1}}}
\newcommand{\strand}[2]{%
  \multicolumn{10}{@{}l@{}}{\rule{0pt}{3.1ex}\textbf{\textcolor{Navy}{#1}}\ \ \sub{#2}}\\[2pt]}

\setlength{\tabcolsep}{3pt}
\renewcommand{\arraystretch}{1.04}

\pagestyle{empty}

\begin{document}

\begin{center}
{\LARGE\bfseries Prior art: who has built which piece}\\[4pt]
{\large\color{Navy} Component matrix for \emph{critical Dicke electrometry}}\\[6pt]
{\footnotesize Extending intracavity Rydberg-EIT microwave sensing (Peng \emph{et al.}, JOSA B \textbf{35}, 2272)
into the superradiant phase transition (Zhang \emph{et al.}, PRL \textbf{110}, 090402).\quad\today}
\end{center}

\vspace{-2mm}
\noindent{\footnotesize
\textbf{Legend:}\ \ \yes\ present \quad \no\ absent \quad
\textcolor{Rust}{\textbf{--}}\,/\,\textcolor{Rust}{\textbf{\checkmark}}\ \emph{the cell that decides the strand.}\ \
\textbf{CRT} = counter-rotating light--matter terms (their presence is what permits a $\mathbb{Z}_2$ superradiant transition).\ \
\textbf{MW = measurand} means the microwave field is the quantity being measured, not merely a drive.}

\vspace{2mm}

{\footnotesize
\begin{tabular}{@{}>{\raggedright\arraybackslash}p{4.8cm} ccc c >{\raggedright\arraybackslash}p{2.9cm} c c c >{\raggedright\arraybackslash}p{4.6cm}@{}}
\toprule
\textbf{Paper} & \textbf{EIT} & \textbf{Cav.} & \textbf{$g\sqrt N$} & \textbf{CRT} &
\textbf{Transition} & \textbf{Ryd.\ int.} & \textbf{MW = meas.} & \textbf{Th/Exp} & \textbf{Headline result} \\
\midrule

\strand{A.\ Cavity + Rydberg electrometry}{everything but the counter-rotating terms $\Rightarrow$ no transition at all}
\rowcolor{StrandA} Peng 2018 \sub{JOSA B 35, 2272}\ \textcolor{Gold}{$\bigstar$} & \yes & \yes & \no\,\sub{weak} & \nokey & \nokey & \no & \yes & Th & $8\times$ lower detectable field; $>\!20\times$ resolution\\
Yang 2020 \sub{JOSA B 37, 1664} & \yes & \yes & \yes\,\sub{coll.\ Rabi} & \nokey & \nokey & \no & \yes & Th & $7\times$ over plain EIT; four-peak transmission\\
\rowcolor{StrandA} Wang 2023 \sub{JOSA B 40, 2604} & \yes & \yes & \yes\,\sub{strong} & \nokey & \nokey & \no & \yes & Th & $196.7\times$ / $26.2\times$; floor $396.5$\,nV\,cm$^{-1}$\\
Jing 2020 \sub{Nat.\ Phys.\ 16, 911} & \yes & \no & \no & \no & \no & \no & \yes\,\sub{superhet} & Exp & $55$\,nV\,cm$^{-1}$Hz$^{-1/2}$\\
\rowcolor{StrandA} Liu 2025 \sub{Chin.\ Phys.\ Lett.\ 42, 053201} & \yes & \no\,\sub{\emph{MW} cav.} & \no & \no & \no & \no & \yes & Exp & 18\,dB power sensitivity\\
arXiv:2502.20792 \sub{cavity superhet} & \yes & \yes & \yes & \nokey & \nokey & \no & \yes\,\sub{superhet} & Exp & 19\,dB sensitivity gain\\

\strand{B.\ Critical Rydberg electrometry}{has criticality \emph{and} sensing --- but the interaction-driven transition, not light--matter}
\rowcolor{StrandB} Ding 2022 \sub{Nat.\ Phys.\ 18, 1447} & \yes & \no & \no & \nokey & \yes\ \sub{NEPT / bistability} & \yes\,\sub{drives it} & \yes & Exp & FI $\times10^3$ vs.\ independent particles; $49$\,nV\,cm$^{-1}$Hz$^{-1/2}$\\
Wang 2025 \sub{arXiv:2502.19761} & \yes & \yes & \no & \nokey & \yes\ \sub{NEPT / bistability} & \yes\,\sub{drives it} & \yes & Exp & classical FI $\times100$ from cavity; $2.6$\,nV\,cm$^{-1}$Hz$^{-1/2}$\\

\strand{C.\ Dicke model / phase structure}{has the transition, the cavity and the interactions --- and senses nothing}
\rowcolor{StrandC} Dimer 2007 / Baden 2014 \sub{PRA 75, 013804 / PRL 113, 020408} & \no & \yes & \yes & \yes\,\sub{Raman} & \yes\ \sub{$\mathbb{Z}_2$ Dicke} & \no & \nokey & Th / Exp & Engineered Dicke QPT via cavity-assisted Raman\\
Zhang 2013 \sub{PRL 110, 090402}\ \textcolor{Gold}{$\bigstar$} & \sub{2-photon} & \yes & \yes & \nokey\,\sub{\textbf{RWA}} & \yes\ \sub{$U(1)$ polariton + SRS} & \yes & \nokey & Th (QMC) & Superradiant solid; first-order boundaries\\
\rowcolor{StrandC} arXiv:2503.13949 \sub{anisotropic engineering} & \no & \yes & \yes & \yeskey\,\sub{tunable $0\!\to\!\infty$} & \yes & \yes & \nokey & Th & Periodic-driving recipe for the anisotropic Dicke model\\
arXiv:2511.22230 \sub{Dicke--Ising QMC} & \no & \yes & \yes & \yes & \yes\ \sub{SR / SRS / Solid-$\tfrac12$} & \yes & \nokey\,\sub{MW = \emph{drive}} & Th (QMC) & Orders of all transitions; blockade suppresses cavity occupation\\
\rowcolor{StrandC} Garbe 20/22, Rams 18, Gietka 21/22, Gyhm 25 & \no & \yes & \yes & \yes & \yes & \no & \no\,\sub{generic param.} & Th & Critical-metrology scaling --- and its no-gos\\

\rowcolor{StrandC} arXiv:1611.00797 \sub{steady-state SR with Rydberg polaritons} & \no & \yes & \yes & \no & \yes\ \sub{lasing threshold} & \yes & \nokey & Th & Rydberg nonlinearity strongly modifies the superradiance threshold\\

\strand{D.\ This proposal}{the same protocol as strand B, with a \emph{light--matter} transition in place of the interaction-driven one}
\rowcolor{Ours} \textbf{Route I} --- $U(1)$ lasing \sub{after Zhang 2013} & \yes & \yes & \yes & \no\,\sub{+ repump} & \yeskey\ \sub{$U(1)$ lasing} & \sub{optional} & \yeskey\,\sub{\textbf{sets} $G^2_{\rm th}$} & Th & cavity drift rejected $101\times$; $0.070$\,nV\,cm$^{-1}$Hz$^{-1/2}$ at $N{=}10^6$\\
\rowcolor{Ours} \textbf{Route II} --- $\mathbb{Z}_2$ Dicke \sub{after Dimer/Baden} & \yes & \yes & \yes & \yeskey & \yeskey\ \sub{$\mathbb{Z}_2$ Dicke} & \sub{optional} & \yeskey\,\sub{\textbf{sets} $\lambda_c$} & Th & $\varepsilon^{-1}$ gain; $-10.3$\,dB self-squeezing\\
\bottomrule
\end{tabular}}

\vspace{1mm}
\noindent{\footnotesize \textcolor{Gold}{$\bigstar$} = the two seed papers this project starts from.\hfill Full discussion in \texttt{dicke\_rydberg\_sensing.pdf}, \S6.\ \ Reading notes and references overleaf.}

\newpage

\section*{\color{Navy} Reading the matrix}

Every individual column is occupied, and each strand is exactly \emph{one column} short of the proposal.

\begin{itemize}[leftmargin=1.4em,itemsep=1pt]
  \item \textbf{Strand A} lacks counter-rotating terms, so there is no transition at all. It stops at a
        collective Rabi splitting read out spectroscopically.
  \item \textbf{Strand B} has criticality and sensing, but its critical point is set by atomic density and
        interaction strength: it drifts with temperature and stray fields, it is hysteretic (the discriminator
        depends on scan direction), and it cannot be tuned quickly.
  \item \textbf{Strand C} has the light--matter transitions and does no sensing. Two near-misses:
        arXiv:2511.22230, where microwave fields are \emph{in} the setup but as the drive that engineers the
        model rather than the measured quantity; and arXiv:1611.00797, which puts a Rydberg medium at a
        superradiance threshold in a cavity, but to generate nonclassical light.
  \item \textbf{Strands B and D are the same phase diagram.} Turning up the Rydberg interaction in Route~I
        folds its threshold over into optical bistability --- i.e.\ back into strand B. The proposal is a move
        between corners of one model, not a competing claim.
\end{itemize}

\begin{tcolorbox}[colback=Ours,colframe=Gold,boxrule=0.8pt,arc=2pt,left=8pt,right=8pt,top=6pt,bottom=6pt]
\textbf{The empty cell that defines the project} is the intersection of the last two columns ---
a \emph{light--matter} transition whose phase boundary \emph{is} the measurand. Both flavours
depend on the same dressed splitting, so both are sensing curves:
\[
  \tilde\omega_0=\delta+\Omega_\mu/2\,,\qquad \Omega_\mu=dE_\mu/\hbar\,;
  \qquad
  G^2_{\rm th}=\kappa\Gamma\!\left[1+\tfrac{(\omega_{\rm cav}-\tilde\omega_0)^2}{(\kappa+\Gamma)^2}\right]
  \ \ (U(1)),
  \qquad
  \lambda_c=\tfrac12\sqrt{\omega\,\tilde\omega_0}\ \ (\mathbb{Z}_2).
\]
\end{tcolorbox}

\vspace{-1mm}
\subsection*{\color{Navy} Two cells worth staring at}

\begin{itemize}[leftmargin=1.4em,itemsep=2pt]
  \item \textbf{Zhang 2013 $\rightarrow$ \nokey\,(RWA).} Their coupling
        $(g/\sqrt N)\sum_i(b_i^\dagger a_i\psi+\mathrm{H.c.})$ conserves excitation number --- which is why their
        phase diagram is drawn against a chemical potential. Their ``generalised Dicke model'' is therefore of
        Tavis--Cummings type and their superradiance is a $U(1)$ polariton condensation, not the $\mathbb{Z}_2$
        transition this proposal needs. \emph{The seed paper is one column short of the required model, so the
        extension is real rather than a re-derivation.}
  \item \textbf{arXiv:2503.13949 $\rightarrow$ \yeskey\,(tunable $0\!\to\!\infty$).} The enabling technology,
        already published --- and it fills exactly the cell Zhang 2013 leaves empty, for the Rydberg platform
        specifically.
\end{itemize}

\noindent The same RWA restriction runs through the whole of strand A. \textbf{Counter-rotating terms are the
single ingredient separating this proposal from that entire lineage.}

\vspace{1mm}
\subsection*{\color{Navy} References}

{\scriptsize
\begin{multicols}{2}
\begin{itemize}[leftmargin=1.1em,itemsep=0pt,parsep=0pt,label={}]
\item Y.~Peng \emph{et al.}, \emph{Cavity-enhanced microwave electric field measurement using Rydberg atoms},
      J.~Opt.~Soc.~Am.~B \textbf{35}, 2272 (2018). \href{https://doi.org/10.1364/JOSAB.35.002272}{doi}
\item A.~Yang \emph{et al.}, \emph{Enhanced measurement of microwave electric fields with collective Rabi
      splitting}, J.~Opt.~Soc.~Am.~B \textbf{37}, 1664 (2020).
\item Y.~Wang \emph{et al.}, \emph{Rydberg-atom-based measurements of microwave electric fields with cavity QED},
      J.~Opt.~Soc.~Am.~B \textbf{40}, 2604 (2023).
\item M.~Jing \emph{et al.}, \emph{Atomic superheterodyne receiver based on microwave-dressed Rydberg
      spectroscopy}, Nat.~Phys.~\textbf{16}, 911 (2020). \href{https://arxiv.org/abs/1902.11063}{arXiv:1902.11063}
\item \emph{Cavity-enhanced Rydberg atom microwave receiver}, Chin.~Phys.~Lett.~\textbf{42}, 053201 (2025).
      \href{https://arxiv.org/abs/2404.06915}{arXiv:2404.06915}
\item \emph{Cavity-enhanced Rydberg atomic superheterodyne receiver}.
      \href{https://arxiv.org/abs/2502.20792}{arXiv:2502.20792}
\item D.-S.~Ding, Z.-K.~Liu, B.-S.~Shi, G.-C.~Guo, K.~M\o{}lmer, C.~S.~Adams, \emph{Enhanced metrology at the
      critical point of a many-body Rydberg atomic system}, Nat.~Phys.~\textbf{18}, 1447 (2022).
      \href{https://arxiv.org/abs/2207.11947}{arXiv:2207.11947}
\item Q.~Wang \emph{et al.}, \emph{High-precision measurement of microwave electric field by cavity-enhanced
      critical behavior\ldots}, Sci.~China Phys.~Mech.~Astron.~(2025).
      \href{https://arxiv.org/abs/2502.19761}{arXiv:2502.19761}
\item F.~Dimer, B.~Estienne, A.~S.~Parkins, H.~J.~Carmichael, Phys.~Rev.~A \textbf{75}, 013804 (2007).
\item M.~P.~Baden \emph{et al.}, \emph{Realization of the Dicke model using cavity-assisted Raman transitions},
      Phys.~Rev.~Lett.~\textbf{113}, 020408 (2014); erratum \textbf{118}, 199901 (2017).
\item X.-F.~Zhang, Q.~Sun, Y.-C.~Wen, W.-M.~Liu, S.~Eggert, A.-C.~Ji, \emph{Rydberg polaritons in a cavity: a
      superradiant solid}, Phys.~Rev.~Lett.~\textbf{110}, 090402 (2013).
      \href{https://arxiv.org/abs/1207.4238}{arXiv:1207.4238}
\item \emph{Engineering anisotropic Dicke model with dipole-dipole interaction for Rydberg atom arrays in
      cavity}. \href{https://arxiv.org/abs/2503.13949}{arXiv:2503.13949}
\item \emph{Quantum phase transitions of the anisotropic Dicke--Ising model in driven Rydberg arrays},
      Phys.~Rev.~A (2026). \href{https://arxiv.org/abs/2511.22230}{arXiv:2511.22230}
\item L.~Garbe, M.~Bina, A.~Keller, M.~G.~A.~Paris, S.~Felicetti, Phys.~Rev.~Lett.~\textbf{124}, 120504 (2020).
\item L.~Garbe, O.~Abah, S.~Felicetti, R.~Puebla, Quantum Sci.~Technol.~\textbf{7}, 035010 (2022).
      \href{https://arxiv.org/abs/2110.04144}{arXiv:2110.04144}
\item M.~M.~Rams \emph{et al.}, \emph{At the limits of criticality-based quantum metrology}, Phys.~Rev.~X
      \textbf{8}, 021022 (2018). \href{https://arxiv.org/abs/1702.05660}{arXiv:1702.05660}
\item J.-Y.~Gyhm, H.~Kwon, M.-J.~Hwang, \emph{Fundamental scaling limit in critical quantum metrology}.
      \href{https://arxiv.org/abs/2506.19003}{arXiv:2506.19003}
\end{itemize}
\end{multicols}}

\end{document}
